\subsubsection{Tiling Security Stuff (Notes)}

\begin{itemize}
    \item self severance
    \item self sev + children
    \item sev child
\end{itemize}

\begin{itemize}
    \item floor((v-1)/2) = just less than half branches of root
    \item root + floor((v-1)/2) -- 50\% always
    \item
\end{itemize}

------

\begin{itemize}
    \item
\end{itemize}

\begin{forest}
    grow'=east,child anchor=west,
    for descendants={grow'=east,child anchor=west, align=center},
    [tiling security structure
     [analysis method
        [construction
            [tile-work]
            [$\text{Sec}(g|i)$ vs $w_{g|i}$]
            [$r$ and $s$]
            [$S {=} \sum_i w_i$]
        ]
     ]
     [evaluate cases
        [control tile \\ requires parent, align=center
            [don't need other cases]
            [below is stricter anyway]
        ]
        [root $> \frac{1}{2}S$, align=center]
        [$\frac{1}{2}$ (r + c) $> qS$, align=center]
        [$\frac{1}{2}$ (r + c + gc) $> qS$, align=center]
     ]
     [concessions
        [recursive PoR \\ validation, align=center
            [validate: \\ root + c + \\ ancestor tiles, align=center
                [$O(c \log(n))$]
            ]
            [problem: how \\ to count work?, align=center
                [sum of work can't \\ be > actual work]
            ]
        ]
     ]
    ]
\end{forest}

\subsubsection{Tree-Tiling Security}

\label{sec:tiling-security}

\todo{``Simplex Tile'' or ``Simplex-Tile''}

What does it mean for a tiling to be secure?
Well, we have our broad definition from the trilemma: the network is secure against an $O(n)$ attacker.
To answer this more specifically, though, we need to first answer: \emph{What does a doublespend against a tiling look like?}
Currently, we don't know enough to answer this -- we next need answers to questions like \emph{Which PoRs do simplex-chains in a given tile use?} and \emph{How do we govern the difficulties of simplex-chains in different tiles?}

To answer these questions, and then evaluate the security of tree-tiling, we first need to define some key concepts.
After that we will construct a security architecture so that we can answer the above questions.
Finally, we will analyze this architecture to find out if it's secure.

\subsubsubsection{Key Concepts}



\subsubsubsubsection{Local Chain-work, Tile-work, and Unit Tile}

We're familiar with \emph{local chain-work} (that's work directly contributed by a single simplex-chain).

\emph{Tile-work} is the work that is contributed from simplex-chains that are part of a given tile.
\emph{Tile-work} is the sum of the local chain-work of each constituent simplex-chain.

We could talk about tile-work directly in terms of \emph{number of hashes}, but this isn't necessary.
Instead, we'll normalize tile-work via a \emph{unit tile}.
If some tile's tile-work is $3\times$ that of the \emph{unit tile}, then we write down that tile-work as $3$ instead of some number of hashes.
This will simplify our analysis.

\subsubsubsubsection{The Core}

\emph{The core} of a tree-tiling is \emph{the root tile and its children}.
It's called \emph{the core} because all nodes in the network verify the work and PoRs of all simplex-chains that are part of the core.
Additionally, the core is where the vast majority of work is done.

\subsubsubsubsection{Recursive \emph{Proof of Reflection}}

\newcommand\cL{\ensuremath{L}\xspace}
\newcommand\cM{\ensuremath{M}\xspace}
\newcommand\cR{\ensuremath{R}\xspace}

Part of our security architecture will be the idea of \emph{recursive PoR}.

Say we have three chains: \cL, \cM, and \cR.
They are doing mutual reflection like this: $\xymatrix@C=1em@M=0em{{\cL} \; \ar@{<->}[r] & \; {\cM} \; \ar@{<->}[r] & \;{\cR}}$.
That is: \cL reflects \cM but not \cR; \cM reflects both \cL and \cR; and \cR reflects \cM but not \cL.

With what we already know, it's clear that \cM is as secure as all three combined.
However, that isn't the case for \cL and \cR.
How can we make \cL and \cR as secure as \cM?
Is this even reasonable?

Let's consider chain \cL -- what happens if \cL's local chain-work is, say, $3\times$ that of chain \cM?
The total chain-work that \cL takes into account is $\nicefrac43$ of its local chain-work, so an attacker would need $\nicefrac23$ worth of \cL's local chain-work to control \cL's history.
This is more than 50\%, but it's within reason that an attacker might achieve the required hash-rate.

Case 1: If \cR's local chain-work is large (say, $3\times$ that of \cM), then \cM is almost twice as hard to attack as \cL.\footnote{
    \cM is $\nicefrac74\times$ harder to attack than \cL.
}
This is a problematic situation: it means that an attacker might be able to perform a doublespend against \cL even though they could not perform one against \cM.
Therefore, this case would not be $O(n)$ secure.

Case 2: However, if \cR's local chain-work is small (say, $\nicefrac13$ that of \cM), then \cR is not particularly significant.
Thus, \cM is not substantially harder to attack than \cL, and an attacker could attack both \cL and \cM simultaneously (optionally, \cR too).
But, \cR \emph{is} substantially easier to attack than either \cM or \cL.
The total chain-work that \cR takes into account is only $\nicefrac4{13}$ that of the whole network.
So, this case is \emph{also} not $O(n)$ secure, even though \cL \emph{is} $O(n)$ secure.

These two problem cases are what recursive PoR solves.

Notice that in both cases, \cM is $O(n)$ secure.
Given this, chains \cL and \cR should be able to use the same work that chain \cM does to reach the same level of security.

The first problem is that \cL and \cR cannot use the same \textsc{WeightOf} algorithm as before (although, \cM can).
Particularly, \cL and \cR's \textsc{WeightOf} must count: not only their local chain-work and that contributed reflected \cM blocks, but also some additional weight based on reflections from \emph{other} chains.
Each \cM block might record the total chain-work that it adds to \cM (including PoRs), but \cL and \cR's new \textsc{WeightOf} can't use this value directly.
Firstly, that value includes past local chain-work that reflected \cM, but more importantly: we don't want to force \cL and \cR to verify \emph{every} PoR that contributes to \cM, just the ones that matter.

There is a trivial solution to this problem: require \cL and \cR nodes to download the headers of \emph{all} relevant chains, verify their PoRs, and track the origin of those PoRs.
Nodes must track the origin of work contributed via PoR to avoid double-counting work.
An easy example of this is: \cL nodes should not count PoRs where an \cL block is the reflecting block; that work is already directly included in \cL's local chain-work.
Another example is where there are \emph{multiple paths} that depend on the same chain.
This leads us into the second problem.

\subsubsubsubsection{Recursive PoR with Multiple Paths}

\newcommand\cMone{\ensuremath{M_1\xspace}}
\newcommand\cMtwo{\ensuremath{M_2\xspace}}

Say that we now have four chains: \cL, \cMone, \cMtwo, and \cR doing mutual PoR like so:
\begin{equation*}
    \xymatrix@C=2em@R=-.15em{
        & {\cMone} & \\
        {\cL} \ar@{<->}[ur] & & {\cR} \ar@{<->}[ul] \\
        & {\cMtwo} \ar@{<->}[ul] \ar@{<->}[ur] &
    }
\end{equation*}

\todo{continue}

\subsubsubsection{Security Architecture}

\subsubsubsubsection{Inter-tile Security Relationship}

$r$

\begin{forest}
    [ITSR
        [govern via RAA + DAA]
        [ratio of work based on $d$
            [$r$ \textrightarrow work drops by $1/r$ per layer]
        ]
        [root tile exception]
    ]
\end{forest}

\subsubsubsubsection{Total Work Across a Tiling}

$S$

\begin{itemize}
    \item Except for the root tile, each tile contributes $\frac1{r^d}$ work to the tiling.
    \item At each depth, there are $u^d$ tiles.
    \item Root tile contributes $M$ units of tile-work.
\end{itemize}

\begin{align}
    S &= M + \frac{u}{r} + \frac{u^2}{r^2} + \cdots && \text{Sum \emph{all} chain-work}
    \nonumber \\
    S - (M-1) &= 1 + \frac{u}{r} + \frac{u^2}{r^2} + \cdots
    \nonumber \\
    \intertext{This is a geometric series, and since $\frac{u}{r} < 1$:}
    S - (M-1) &= \frac{1}{1 - \frac{u}{r}} && \text{Sum of geometric series}
    \nonumber \\
    S - (M-1) &= \frac{r}{r - u} && \times\frac{r}r
    \nonumber \\
    S &= M-1 + \frac{r}{r - u}
    \label{eq:tree-tiling-total-work}
\end{align}

\subsubsubsubsection{Doublespend Requirements}

\subsubsubsection{Analysis Method}

How can we analyze the security of tree-tilings?


\subsubsubsection{Assumptions}

\begin{itemize}
    \item security of tile works like that of a simplex
    \item able to apply logic recursively
\end{itemize}

\subsubsubsection{Additional Imposed Requirements}

\subsubsubsection{Security Criteria for Specific Cases}

\subsubsubsection{Security Evaluation of Specific Cases}

\begin{itemize}
    \item $r>u$
    \item \dots
\end{itemize}

\subsubsubsubsection{The Core is (Almost) as Secure as the Network}

$M \ge 1$, $u > 1$, $r > u$, and $0 \le q \le 0.5$.

\begin{align}
    \frac12 (M + \frac{u}r) &> qS
    \nonumber \\
    \frac12 (Mr + u) &> rq(M-1+\frac{r}{r-u}) && \times r \text{ and substitute } S \text{ via \autoref{eq:tree-tiling-total-work}}
    \nonumber \\
    (Mr + u)(r-u) &> 2rq((M-1)(r-u)+r) && \times 2(r-u)
    \label{eq:tree-tiling-safe-q}
    % https://www.wolframalpha.com/input?i=%28My+%2B+x%29%28y-x%29+%3E+2yq%28%28M-1%29%28y-x%29%2By%29%2C+0+%3C%3D+q+%3C%3D+0.5
    % note for above: y=r; x=u
    % 2Mr &> |u| \sqrt{\frac{2M^2q-M^2-4Mq-2M+2q-1}{2q-1}}+u(M-1) \\
    % 2Mr &> |u| \sqrt{\frac{M^2(1 - 2q) - 2M(2q + 1) + (2q - 1)}{2q-1}}+u(M-1) \\
\end{align}

There is no elegant form of \autoref{eq:tree-tiling-safe-q}, but this form is, at least, easy to evaluate and graph.

\subsubsubsection{Scaling Complexity}

\subsubsubsubsection{Complexity of the Core}

\begin{forest}
    rectangle, draw=black, minimum width=6em, minimum height=1.5em,
    for descendants={rectangle, draw=black, minimum width=3em, minimum height=1.5em}
    [
        0,
        [0|0
            [0|0|0]
            [0|0|1]
        ]
        [0|1
            [0|1|0]
            [0|1|1]
        ]
    ]
\end{forest}

chains:

\begin{align*}
    N_\text{chains} &= N_{1,d=0} + 2N_{1,d=1} + 4N_{1,d=2} \\
    &= \frac{N_1}{4} + 2\frac{N_1}{4} + 4\frac{N_1}{4} \\
    &= \frac74 N_1 = O(c) \\
    \intertext{but if the core is just the root + children:}
    N_\text{chains} &= \frac34 N_1
\end{align*}

refls

\providecommand\parens[1]{%
{\left( #1 \right)}%
}
\providecommand\NOneSq{{{N_1}^2}}

\begin{align*}
    N_\text{refls} &= \left(\frac{N_1}{4}\right)^2 + 6\left( \frac{N_1}{4} \right)^2
                    + 2\left(\frac{N_1}{4}\cdot \frac{N_1}{4}\right)
                    + 4\left(\frac{N_1}{4}\right)^2 \\
    &= \NOneSq\parens{\frac{1}{16} + \frac{6}{16} + \frac{2}{16} + \frac{4}{16}} \\
    &= \frac{13}{16}\NOneSq = O(c^2) \\
    \intertext{but if the core is just the root + children:}
    N_\text{refls} &= \parens{\frac{N_1}{4}}^2 + 2\parens{\frac{N_1}{4}}^2 + 2\parens{\frac{N_1}{4}\cdot \frac{N_1}{4}} \\
    &= \frac{5}{16}\NOneSq
\end{align*}

for a leaf tile, each layer adds $2\parens{\frac{N_1}{4}}^2$ more reflections, so we get $\ceil{8 (1 - \frac5{16})}-1 = 5$ layers for free.

\begin{align*}
    2^2 - 1 &\to 2^{2+5} - 1 \\
    &\implies \frac{2^7 - 1}{3}\times \text{ the capacity.} \\
    &\approx 42\times
\end{align*}

\begin{equation*}
    a
\end{equation*}

\subsubsubsubsection{Security Burden Complexity}

each tile has a local $O(c)$ load for PoR verification

each tile is $h \lesssim \log_2(n)$ away from the root

all tiles between that tile and the root need to have PoRs verified, plus the root and its surrounding tiles, and the local tile's neighbors.

root + children + gc + in-between + local + children

$R = \text{root} = 2$

that's $((v-1)^2 + (v-1) + R) + \max(0, h - 3) + (1 + (v-1))$

\begin{align*}
    & & &\quad\:\: ((v-1)^2 + (v-1) + R) + \max(0, h - 3) + (1 + (v-1)) \\
    & & &= (v^2 - 2v - 1 + v - 1 + 2) + \max(0, h - 3) + v \\
    & & &= v^2 + \max(0, h-3) \\
    & & &= O(v^2) + O(\max(0, h-3)) \\
    & & &= O(v^2) + O(h) \\
    & & &= O(v^2) + O(\log(n)) \\
    & & &= O(v^2 + \log(n)) \\
    \text{since } v^2 \text{ is constant:} & & &= O(\log(n)) \\
    % \therefore \text{ (assuming +PoRs) total load is:} & & & O(c\log(n)) \\
    \therefore \text{ total load is:} & & & O(c^2\log(n)) \\
\end{align*}

\subsubsubsubsection{Capacity Complexity}

\begin{center}
\begin{forest}
    [Root
        [c1
            [c11]
            [c12]
        ]
        [c2
            [c21]
            [c22]
        ]
    ]
\end{forest}
\end{center}

\begin{align*}
    N_\text{tiles} &= \frac{u^{h+1} - 1}{u-1} \\
    &= O(u^h) = O(n)
\end{align*}

\begin{align*}
    2^{44} &> r^{h} \\
    44 &> h \log_2(r) \\
    \frac{44}{\log_2(r)} &> h \\
    \intertext{$r' \ge r$, so:}
    \frac{44}{\log_2(r)} \ge \frac{44}{\log_2(r')} &> h \\
    \frac{44}{\log_2(5)} &> h \\
    \sim 19 &> h \\
    \implies N_\text{tiles} &\lesssim 2^{19} \approx 5\cdot 10^{5} \\
    \intertext{if 96 bits are available:}
    \sim 41 &> h \\
    \implies N_\text{tiles} &\lesssim 2^{41} \approx 2\cdot 10^{12} \\
\end{align*}

\subsubsection{Asymmetry: Capacity and Security}

% https://www.desmos.com/calculator/7trq5aehnm
% https://www.desmos.com/calculator/n17w3560ig

\todo{why do $n$-ary tree-tilings?}
